\section{Sum of Subarray Minimums}
\label{sec:sq-sum-subarray-minimums}

\problemstatement{Given an array $\textit{arr}$ of $n$ integers, consider
every one of its $\binom{n+1}{2}$ contiguous subarrays. Compute the sum,
over all of them, of each subarray's minimum element (the result can be
large, so it is typically returned modulo $10^9+7$).}

\begin{patternbox}
\emph{``Sum, over all subarrays, of each subarray's minimum (or
maximum)''} is a strong signal to flip the entire computation around:
instead of asking, for each subarray, ``what is its minimum,'' ask, for
each \textbf{element}, ``in how many subarrays is \emph{this} element the
minimum.'' This contribution-counting reformulation, combined with a
monotonic stack to compute exactly that count for every element in
$O(n)$ total, avoids ever enumerating the $O(n^2)$ subarrays directly.
\end{patternbox}

\subsection*{Understanding the problem carefully}

For element $\textit{arr}[i]$ to be the minimum of some subarray, that
subarray must be entirely contained within the maximal range around $i$
over which $\textit{arr}[i]$ remains the smallest value -- bounded on the
left by the nearest strictly smaller element, and on the right by the
nearest smaller-or-equal element (using a strict boundary on one side and
a non-strict boundary on the other is essential to avoid double-counting
subarrays when duplicate values are present, by consistently attributing
each subarray's minimum to a single, well-defined occurrence among any
tied values).

\begin{ideabox}[{Count Subarrays Where $\textit{arr}[i]$ Is the Minimum}]
For each $i$, let $\textit{left}[i]$ be the number of valid choices for a
subarray's left endpoint such that every element from that endpoint up to
$i$ is $\ge \textit{arr}[i]$ -- this equals $i - \textit{PLE}(i)$, where
$\textit{PLE}(i)$ is the index of the nearest element to the left that is
\emph{strictly less} than $\textit{arr}[i]$ (or $-1$ if none). Let
$\textit{right}[i]$ be the number of valid choices for the right endpoint
such that every element from $i$ to that endpoint is $\ge
\textit{arr}[i]$ -- this equals $\textit{NLE}(i) - i$, where
$\textit{NLE}(i)$ is the nearest element to the right that is \emph{less
than or equal to} $\textit{arr}[i]$ (or $n$ if none). Then
$\textit{arr}[i]$ is the minimum of exactly $\textit{left}[i] \times
\textit{right}[i]$ subarrays (any combination of a valid left endpoint and
a valid right endpoint).
\end{ideabox}

\subsection*{Why the strict/non-strict boundary asymmetry avoids double-counting}

If both boundaries used the same (either both strict, or both
non-strict) comparison, a subarray containing two equal minimum values
would be counted once for \emph{each} occurrence of that value, inflating
the sum. Using strictly-less on the left and less-or-equal on the right
ensures that, among several equal candidate minimum positions within a
subarray, only the \textbf{leftmost} one ``claims'' that subarray as
contributing to its count: the leftmost equal occurrence's right boundary
search will stop at the very next equal (or smaller) value, while any
later equal occurrence's left boundary search will pass straight over
the earlier equal value (since it is not \emph{strictly} less) and only
stop at a genuinely smaller one further left -- so the two occurrences
never both count the same subarray.

\subsection*{Algorithm}

\begin{enumerate}
  \item Compute $\textit{left}[i]$ for every $i$ via a monotonic
        increasing stack scanned left to right: pop while the stack's top
        value $\ge \textit{arr}[i]$; $\textit{left}[i] \gets i - (\text{
        stack empty} ? -1 : \text{top index})$; push $i$.
  \item Compute $\textit{right}[i]$ for every $i$ via a monotonic stack
        scanned right to left: pop while the stack's top value $>
        \textit{arr}[i]$; $\textit{right}[i] \gets (\text{stack empty} ?
        n : \text{top index}) - i$; push $i$.
  \item Sum $\textit{arr}[i] \times \textit{left}[i] \times
        \textit{right}[i]$ over all $i$, modulo $10^9+7$.
\end{enumerate}

\subsection*{Dry run}

\dryrun{Take $\textit{arr} = [3,1,2,4]$.}

\textbf{Left} (pop while top $\ge \textit{arr}[i]$): $\textit{left} =
[1,2,1,1]$ (computed via the standard PLE scan, e.g.\ at $i=2$, value
$2$: the stack's top is index $1$, value $1 < 2$, so no pop;
$\textit{left}[2]=2-1=1$).

\textbf{Right} (pop while top $> \textit{arr}[i]$): $\textit{right} =
[1,3,2,1]$ (e.g.\ at $i=1$, value $1$: scanning right to left, both
index $2$ (value $2$) and index $3$... after appropriate pops the stack
empties, giving $\textit{right}[1]=4-1=3$).

Contributions: $3 \times 1 \times 1=3$; $1\times2\times3=6$;
$2\times1\times2=4$; $4\times1\times1=4$. Total $=3+6+4+4=17$.

\subsection*{C++ implementation}

\begin{lstlisting}
#include <bits/stdc++.h>
using namespace std;
const int MOD = 1e9 + 7;

int sumSubarrayMins(vector<int>& arr) {
    int n = arr.size();
    vector<int> left(n), right(n);
    stack<int> st;

    for (int i = 0; i < n; i++) {
        while (!st.empty() && arr[st.top()] >= arr[i]) st.pop();
        left[i] = i - (st.empty() ? -1 : st.top());
        st.push(i);
    }

    while (!st.empty()) st.pop();

    for (int i = n - 1; i >= 0; i--) {
        while (!st.empty() && arr[st.top()] > arr[i]) st.pop();
        right[i] = (st.empty() ? n : st.top()) - i;
        st.push(i);
    }

    long long total = 0;
    for (int i = 0; i < n; i++) {
        total = (total + (long long)arr[i] * left[i] % MOD * right[i]) % MOD;
    }
    return total;
}

int main() {
    vector<int> arr = {3, 1, 2, 4};
    cout << "Sum of subarray minimums: " << sumSubarrayMins(arr) << endl;
    return 0;
}
\end{lstlisting}

\begin{complexitybox}
\textbf{Time Complexity.} Two monotonic-stack passes, each $O(n)$, giving
\[
  O(n).
\]
\textbf{Space Complexity.} $O(n)$ for the stacks and the
\textit{left}/\textit{right} arrays.
\end{complexitybox}

\begin{takeawaybox}
``Sum/count over all subarrays of [some per-subarray extremum]'' is
solved not by enumerating subarrays, but by computing, for each element,
\emph{how many} subarrays it is the extremum of -- a count derived from
the nearest strictly-more-extreme boundaries on each side, found with two
monotonic-stack passes. This contribution-counting reformulation is one
of the most powerful general techniques to recognize in this chapter.
\end{takeawaybox}
